\documentclass[conference]{IEEEtran}

\usepackage{xcolor}
\usepackage[utf8]{inputenc}
\usepackage[T1]{fontenc}
\usepackage[brazilian]{babel}
\usepackage{array}
\usepackage[table]{xcolor}
\usepackage{makecell}
\newcommand{\cmark}{\textbf{X}}
\usepackage{graphicx}
\usepackage{booktabs}
\usepackage{amsmath}
\usepackage[hyphens]{url}
\usepackage[colorlinks=true,allcolors=black]{hyperref}
\usepackage[inline]{enumitem}
\usepackage{graphicx}
\usepackage{float}
\usepackage{tabularx}
\usepackage{array}
\usepackage{svg}
\usepackage{threeparttable}


\newcommand{\gaplit}[1]{%
\par\noindent\colorbox{yellow}{\parbox{0.94\linewidth}{\footnotesize\itshape O que ainda precisamos pesquisar: #1}}%
\par\smallskip%
}

\title{Seleção Dinâmica de Classificadores para \\ Detecção Multimodal de Sarcasmo}

\author{
\IEEEauthorblockA{Experiência Criativa: Projeto Transformador II\\
Bacharelado em Ciência da Computação -- PUCPR\\[2ex]
Turma 8B -- Equipe 7\\
Ana Flávia Martins dos Santos, Isabella Vanderlinde Berkembrock,\\
Michele Cristina Otta, Yejin Chung\\
{\tt\small \{s.flavia2, isabella.berkembrock, michele.otta, yejin.chung\}@pucpr.edu.br}
}
}

\date{today}

\sloppy
\begin{document}

\maketitle

\section*{Resumo}

O sarcasmo é um fenômeno complexo, difícil até para humanos, pois se manifesta por hábitos individuais de expressão e diferentes tipos, com múltiplas modalidades e formas distintas de incongruência. Este trabalho propõe aplicar a Seleção Dinâmica (DS) de classificadores à detecção de sarcasmo, usando a incongruência intermodal como critério explícito de competência. Com a base MUStARD++, foram avaliadas as estratégias OLA, KNORA-E e KNORA-U, comparadas a classificadores estáticos e a uma arquitetura Soft-MoE. Nos 10 \textit{splits} por falante, os classificadores estáticos (SVM e voto dos três modelos) obtiveram o maior F1-macro médio (80.1). A melhor DS foi o KNORA-U, com F1-macro médio entre 79.2 e 79.3 nas suas variantes, e o Soft-MoE obteve 78.3. O teste de Friedman indicou diferença global entre os métodos, mas nenhum par diferiu no pós-teste de Nemenyi, e o roteamento por incongruência não alterou os resultados de forma significativa. Ainda assim, o oráculo atingiu 91.0 de acurácia, o que evidencia uma margem para melhorar a seleção.

\section{Introdução}
\label{sec:intro}

Sarcasmo é um fenômeno em que o significado pretendido diverge do literal, sendo uma forma de expressão de emoções humanas complexas. Sua detecção é relevante para serviços com interação humana, como \textit{chatbots} e \textit{e-commerce}, visto que interpretações equivocadas podem levar a respostas imprecisas, consolidando-se como um desafio para as áreas de análise de sentimentos e processamento de linguagem natural~\cite{chang2026}.

O sarcasmo pode ser classificado em quatro tipos: Proposicional (PRO), que exige informação contextual para determinar se a proposição contradiz a intenção real do falante; Incorporado (EMB), em que a incongruência está contida dentro de uma única sentença, o próprio texto é suficiente para detectá-la; Prefixado (\textit{Like-Phrases}, LIK), utiliza marcadores como "como se" ou "tipo" para expressar a incongruência do argumento (exemplo: "como se você se importasse"); e Ilocutório (ILL), que envolve discrepâncias emocionais entre diferentes modalidades, nas quais o texto normalmente indica uma atitude oposta à capturada por áudio ou vídeo (como exemplo, revirar os olhos ao dizer que está tudo bem). Enquanto os tipos EMB, LIK e PRO baseiam-se principalmente no conteúdo textual — dependendo, no caso de LIK e PRO, de raciocínio pragmático ou contextual para uma interpretação correta —, o sarcasmo ilocutório (ILL) é considerado o mais desafiador por envolver elementos não verbais (\textit{non-speech}), como linguagem corporal e expressões faciais~\cite{chang2026, ray2022multimodal, bhosale2023mustard++}.

Dessa forma, tratar o sarcasmo como um fenômeno singular e monolítico, desconsiderando suas especificidades categóricas e baseando-se unicamente em pistas textuais, compromete severamente os sistemas de análise de sentimento (especialmente os fundamentados em polaridade lexical). Como a incongruência que caracteriza o sarcasmo se manifesta de formas distintas em cada tipo, as pistas multimodais tornam-se complementares, visto que nenhuma modalidade isolada consegue capturá-lo de maneira consistente em todos os cenários~\cite{ray2022multimodal, castro2019towards}. Além disso, essa evidência multimodal não se distribui de maneira uniforme entre as instâncias: em categorias como EMB ou PRO, a informação textual pode ser suficiente ou redundante entre modalidades, enquanto no sarcasmo ILL a evidência emerge predominantemente do desacordo entre os canais de áudio, vídeo e texto~\cite{ray2022multimodal, yu2024mmoe}. Essa variação — combinada aos hábitos individuais de expressão descritos por Du et al.~\cite{du2022effective} — explica por que modelos treinados sem a sensibilidade a essas nuances generalizam mal entre diferentes \textit{datasets}~\cite{jang2024generalizable}.

Apesar disso, os esquemas tradicionais de fusão na detecção de sarcasmo multimodal são estáticos: aplicam o mesmo operador fixo a todas as instâncias~\cite{farabi2024survey}. A Seleção Dinâmica de Classificadores (\textit{Dynamic Selection}, DS) trata esse problema condicionando a predição final a um mecanismo por instância e não a um operador fixado a priori; na formulação clássica da literatura, a competência de cada classificador é estimada por sua vizinhança no espaço de \textit{features}, mas esse é apenas um critério possível entre outros~\cite{britto2014dynamic, cruz2018dynamic}.  Neste trabalho, propõe-se investigar explicitamente a incongruência multimodal como um desses sinais.

Para cada instância, foi proposto estimar a emoção associada às modalidades textual, contextual, acústica e visual e, a partir dela, a polaridade de cada canal (positiva, negativa ou neutra). A comparação dessas polaridades indica se as modalidades concordam ou divergem, e essa informação é usada para agrupar as instâncias. A hipótese é que instâncias com padrões de incongruência distintos favoreçam classificadores distintos, de modo que restringir a estimativa de competência a instâncias do mesmo grupo permita selecionar, por instância, o classificador ou conjunto de classificadores mais adequado.

Arquiteturas de Mistura de Especialistas (\textit{Mixture of Experts}, MoE) constituem uma família relacionada, já aplicada ao sarcasmo multimodal com intuição semelhante de especialização por instância~\cite{yu2024mmoe, zhao2025beyond}, com a diferença de que seus \textit{gates} são tipicamente aprendidos ponta a ponta, enquanto a proposta deste trabalho investiga um sinal explícito e interpretável de incongruência entre modalidades como informação para a seleção dinâmica.

A pergunta que guiou o trabalho foi: \textit{A integração entre Seleção Dinâmica de Classificadores e incongruência intermodal, como critério explícito de competência, melhora a detecção multimodal de sarcasmo em relação a uma rede MoE tradicional?} Para responder adequadamente a pergunta de pesquisa, este trabalho pretende: (i) construir um \textit{pool} heterogêneo de classificadores, cobrindo diferentes combinações de modalidades e famílias de aprendizado; (ii) estimar a valência e a polaridade emocional das modalidades textual, contextual, acústica e visual; (iii) agrupar as instâncias segundo a concordância ou o conflito de polaridade entre modalidades; (iv) investigar o uso desses grupos na definição da região de competência da seleção dinâmica; (v) medir o teto de adaptação por instância por meio da comparação com um oráculo; e (vi) avaliar a proposta contra baselines que permitam isolar tanto a contribuição da seleção dinâmica quanto a contribuição específica dos grupos de incongruência.

A principal contribuição deste trabalho é uma avaliação controlada do uso da incongruência intermodal como critério explícito de competência na seleção dinâmica para a detecção multimodal de sarcasmo, uso que, na literatura consultada, não havia sido explorado. Como restringir a vizinhança altera a estimativa de competência independentemente do critério adotado, cada roteamento por incongruência foi confrontado com controles aleatórios de mesmo tamanho, o que separa o efeito da incongruência do efeito da própria restrição. O resultado foi negativo: a incongruência, tal como foi possível medi-la, não melhorou a seleção. A comparação dos extratores de emoção com a anotação da base sugere uma explicação, já que os grupos refletiram mais a ausência de emoção detectada do que o desacordo entre as modalidades.

% Se sobra paragrafo de organização do artigo;

\section{Trabalhos Relacionados}
\label{sec:related_work}

% REVISAR E MELHORAR ***************************************************
% sempre que possível, recomenda-se mencionar os conjuntos de dados empregados, as métricas adotadas e os resultados reportados nos estudos revisados
% Para cada artigo ou grupo de estudos, sintetizem a técnica, metodologia ou arquitetura proposta, ressaltando as principais contribuições, diferenciais e limitações observadas.

Ray et al. propuseram uma abordagem multimodal para a detecção de sarcasmo nas bases MUStARD e MUStARD++. A proposta envolveu uma arquitetura de \textit{collaborative gating} para fundir diretamente as \textit{features} multimodais, seguida da classificação por camadas totalmente conectadas. Os hiperparâmetros foram selecionados por uma validação cruzada \textit{5-fold} e foram realizados testes com 28 configurações diferentes: independente ou dependente de falante, com ou sem contexto e com 7 combinações entre as modalidades (texto, áudio e visual). A dependência de falante, neste caso, refere-se a passar informação do falante ao modelo. Embora os resultados dependente de falante tenham sido ligeiramente superiores, os autores relataram que uma configuração independente seria mais adequada, já que essa dependência pode enviesar os modelos e reduz a generalização \cite{wang2026bottleneck}. Para o MUStARD++ o maior F1-Score obtido foi de 70.3 nas configurações dependente de falante, com contexto e com fusão das três modalidades \cite{ray2022multimodal}.

Bhosale et al. empregaram uma arquitetura com mecanismo de atenção por \textit{gating} colaborativo sob uma configuração independente de locutor avaliada na base MUStARD++. O protocolo baseou-se em validação cruzada estratificada \textit{5-fold}, com teto de 50 épocas e parada antecipada. O melhor desempenho atingido foi um F1-Score de 73.1 nas configurações com contexto e fusão das três modalidades~\cite{bhosale2023mustard++}.

Para superar as limitações das abordagens de fusão estática tradicionalmente empregadas e capturar a heterogeneidade intermodal, pesquisas incorporaram arquiteturas \textit{Mixture of Experts} na detecção multimodal de sarcasmo. A proposta de Yu et al. \cite{yu2024mmoe} teve como foco lidar com diferentes tipos de interações entre as modalidades de texto e imagem: redundância, unicidade e sinergia. As amostras foram categorizadas na fase de treinamento a partir do grau de concordância entre os modelos unimodais e o rótulo multimodal real. Para a detecção de sarcasmo na base MUStARD, a combinação de Qwen2-1.5B+MMoE obteve uma acurácia de 70.69 e F1-Score de 72.34. Jana et al.~\cite{jana2025midre} propuseram o MiDRE (\textit{Mixture of Dual Reasoning Experts}), que combina um especialista de raciocínio interno, voltado à incongruência entre imagem e texto, e um especialista externo baseado em justificativas geradas por um \textit{Large Vision-Language Model} via \textit{Chain-of-Thought}, ponderados por um \textit{gating} adaptativo. Foi alcançado um F1 de 87.79 na base MMSD2.0. O trabalho superou a baseline, mas operou sobre pares imagem-texto e dependeu de portas aprendidas de ponta a ponta.

% esse cara fez DS pqp >>> Zhao et al \cite{zhao2025beyond} propuseram o MM-MoE, um framework de MoE com \textit{gating} dinâmico global para fusão \textit{cross-modal} adaptativa, visando reduzir a dependência de pistas espúrias, como emojis e hashtags.
Técnicas de Seleção Dinâmica de Classificadores oferecem uma alternativa fundamentada para a adaptação por instância, avaliando a competência local de um conjunto de classificadores na região de competência de cada amostra de teste \cite{britto2014dynamic,cruz2018dynamic}. Entre os diferentes métodos para a seleção de especialistas pode-se mencionar o DS-LA OLA proposto por Woods et al. \cite{woods1997combination}, baseado na acurácia, e métodos K-Nearest-Oracles (KNORA), como o Knora-Eliminate e Knora-Union, propostos por Ko et al. \cite{ko2008dynamic}. Embora amplamente consolidadas no aprendizado de máquina padrão, o uso de estratégias de DS no domínio de detecção multimodal de sarcasmo ainda permanece pouco explorado. 


% FAZERRR: *********************************************************
% Ao final da seção, apresentem uma síntese geral destacando os avanços, limitações e lacunas identificadas na literatura, indicando como o presente trabalho se insere nesse contexto e quais aspectos busca aprimorar ou explorar !!!!!!!!!!!!!!!!!!!!!!!!!!

Em síntese, os estudos revisados ressaltam a importância do contexto para a detecção de sarcasmo e demonstram que arquiteturas MoE já exploram a heterogeneidade intermodal. Contudo, essas abordagens ainda dependem de esquemas de fusão estáticos ou de portas aprendidas de ponta a ponta, o que limita a adaptabilidade e interpretabilidade da decisão para cada instância. Na literatura consultada, ainda há carência de trabalhos com aplicação da DS à detecção multimodal de sarcasmo, e não foi explorado o uso explícito da incongruência intermodal na definição da região de competência. Este trabalho busca preencher essa lacuna ao investigar as estratégias de DS OLA, KNORA-E e KNORA-U, na base MUStARD++, incorporando a incongruência intermodal à definição da região de competência.

\section{Materiais e Métodos}
\label{sec:materiais}

% O método proposto integra representação multimodal, estimação de valência para cálcula da incongruência, formação de uma \textit{pool} heterogênea de classificadores e DS. A partir das modalidades de texto, contexto, áudio e visual, são extraídas representações utilizadas tanto na classificação quanto na identificação de padrões de incongruência emocional. Esses sinais orientam a definição da região de competência utilizada pelos métodos de seleção dinâmica. A Figura~\ref{fig:pipeline} apresenta uma visão geral do procedimento experimental adotado.

A Figura~\ref{fig:pipeline} apresenta uma visão geral do procedimento experimental adotado, que integra representação multimodal, estimação de valência para o cálculo da incongruência, formação de uma \textit{pool} heterogênea de classificadores e DS. Os detalhes de cada etapa são descritos nas seções seguintes.

\begin{figure}[H]
  \centering
  \includegraphics[width=0.99\linewidth]{images/diagrama TCC.png}
  \caption{Exemplificação do procedimento experimental proposto.}
  \label{fig:pipeline}
\end{figure}

\subsection{Base de Dados}

Para a execução do projeto, foi selecionada a base de dados MUStARD++\footnote{Disponível em: \url{https://github.com/cfiltnlp/MUStARD_Plus_Plus}}, uma versão extendida da base MUStARD com adição de rótulos e balanceada em 601 instâncias sarcásticas e 601 não sarcásticas, totalizando 1202 instâncias. Os dados foram retirados de programas de sit-com, como: The Big Bang Theory, Friends, The Golden Girls, Sarcasmaholics Anonymous e The Silicon Valley \cite{ray2022multimodal, castro2019towards}. Os atributos encontrados na base estão detalhados na Tabela \ref{tab:dataset} e a Figura \ref{fig:dataset} ilustra um exemplo de instância.

\input{tables/dataset}

As variáveis de emoção implícita, emoção explícita, valência e excitação (\textit{arousal}) foram mensuradas considerando todas as modalidades combinadas e, junto do tipo de sarcasmo, só possuem valores em amostras alvo do \textit{dataset}.

% Figura de exemplo
\begin{figure}[h] 
    \centering 
    \includegraphics[width=0.99\linewidth]{images/diagrama mustard.png} 
    \caption{Exemplo de instância do \textit{dataset}} 
    \label{fig:dataset} 
\end{figure}


\subsection{Recursos computacionais}
Os experimentos foram conduzidos em uma máquina com processador Intel Core i7-13620H, 16 GB de memória RAM e GPU NVIDIA GeForce RTX 4050 \textit{Laptop} (6 GB). 

\subsection{Representação das Modalidades}
% lembrar de add refs depois************************************

Para cada instância do MUStARD++, foram extraídas representações independentes para as modalidades de texto, contexto, áudio e visual. Essas representações foram posteriormente utilizadas como entrada dos classificadores individuais e de suas diferentes combinações multimodais.

Para a modalidade textual, foi utilizada a transcrição do enunciado-alvo da instância, processado pelo modelo pré-treinado \textit{BART} \cite{lewis2020bart}, sendo utilizada a média dos estados ocultos para produzir um vetor de representação semântica de 768 dimensões. O contexto foi processado separadamente, a partir dos enunciados anteriores associados à cena, seguindo o mesmo procedimento, também resultando em 768 dimensões.

Na modalidade acústica, o áudio correspondente ao enunciado-alvo foi processado pelo modelo pré-treinado \textit{Wav2Vec2} \cite{baevski2020wav2vec}. As representações produzidas ao longo do tempo foram agregadas por meio da média temporal, produzindo um vetor acústico de 768 dimensões para cada instância.

% ref********************************
Para a modalidade visual, foram extraídos quadros do vídeo correspondente a cada instância e, em cada quadro, foi realizada a detecção das faces presentes. As faces detectadas foram processadas por um modelo pré-treinado baseado em \textit{Vision Transformer} (ViT) para reconhecimento de expressões faciais, responsável por atribuir probabilidades a sete categorias emocionais: \textit{angry}, \textit{disgust}, \textit{fear}, \textit{happy}, \textit{neutral}, \textit{sad} e \textit{surprise}. Para cada instância, foram calculadas a média e o desvio-padrão temporal das probabilidades de cada categoria, além do número de rostos detectados e de um indicador de ausência de informação visual. A representação visual resultante possui, portanto, 16 características.

 Antes do treinamento dos classificadores, as representações correspondentes às modalidades selecionadas em cada configuração foram concatenadas, permitindo avaliar separadamente modalidades individuais e diferentes combinações entre texto, contexto, áudio e visual.

% cite{ma2024emotion2vec} 

\subsection{Valência por Modalidade e Grupos de Incongruência}
\label{subsec:incongruencia}

As anotações de valência e de emoção do MUStARD++ descrevem a fala com todas as modalidades em conjunto e, por isso, não permitem observar o desacordo entre canais. A emoção de cada modalidade foi então estimada por classificadores pré-treinados: um DeBERTa-v3-large~\cite{he2023debertav3} ajustado para sete emoções~\cite{tanneru2025deberta}, para texto e contexto; o emotion2vec+ \textit{large}~\cite{ma2024emotion2vec}, para o áudio; e o ViT de expressão facial \texttt{mo-thecreator/vit-Facial-Expression-Recognition}~\cite{dosovitskiy2021vit, motheecreator_vitfer}, para o rosto. As saídas foram mapeadas para as mesmas sete categorias, e a valência de cada modalidade $m$ foi calculada pela Equação~\ref{eq:valencia}, com \textit{neutral} e \textit{surprise} contribuindo com zero.

\begin{equation}
\label{eq:valencia}
V_m = P_m(\textit{happy}) - \sum_{e \in \mathcal{N}} P_m(e) \in [-1, 1]
\end{equation}
em que $\mathcal{N} = \{\textit{angry}, \textit{disgust}, \textit{fear}, \textit{sad}\}$.

A polaridade de cada modalidade foi dada pela emoção predominante: positiva, negativa ou neutra, sendo neutra também a ausência de sinal (por exemplo, vídeo sem rosto detectado). Com ela, foram definidos dois agrupamentos A/B: a instância pertence ao grupo A quando as quatro modalidades possuem a mesma polaridade e ao grupo B caso contrário (compostos por 173 e 1029 instâncias, respectivamente). Como essa regra conta a neutralidade como discordância, definiu-se também o agrupamento por conflito, em que basta um par de modalidades com polaridades opostas (resultando em 330 instâncias com conflito e 872 sem). Em nenhum dos casos a incongruência entra como \textit{feature} dos classificadores: ela apenas restringe a região do DSEL em que a competência é estimada.


\subsection{Pool de Classificadores}

Para a formação da \textit{pool} de classificadores, foram consideradas todas as 15 combinações não vazias das quatro modalidades disponíveis: texto, contexto, áudio e visual, incluindo combinações unimodais, bimodais, trimodais e multimodal completa. Para cada combinação, foram avaliadas três famílias de classificadores: Random Forest (RF), Regressão Logística (LR) e Support Vector Machine (SVM), totalizando 45 candidatos à \textit{pool}.

Os candidatos foram avaliados por validação cruzada estratificada \textit{5-fold}, respeitando o agrupamento por falante (\texttt{StratifiedGroupKFold}). Os mesmos \textit{folds} foram utilizados para todos os candidatos, permitindo uma comparação consistente entre as diferentes configurações. Para cada candidato, realizou-se uma busca em grade (\texttt{GridSearchCV}) para determinar a melhor configuração de hiperparâmetros por F1-macro. A partir das configurações selecionadas, foram geradas predições \textit{Out-of-Fold} (OOF), utilizadas para calcular o \textit{disagreement} entre pares de candidatos, isto é, a proporção de instâncias em que dois classificadores divergem, como medida de diversidade \cite{monteiro2021diversity, kuncheva2003measures}.

A formação da \textit{pool} considerou não apenas o desempenho individual dos candidatos, mas também a diversidade entre eles. Essa consideração parte do princípio de que modelos com desempenhos elevados, mas com comportamentos muito semelhantes, podem produzir erros redundantes, reduzindo o benefício da combinação de especialistas. Assim, buscou-se a construção de uma \textit{pool} heterogênea, composta por especialistas com bom desempenho individual e padrões de erro complementares. A diversidade foi explorada em diferentes dimensões: no espaço de classificadores, por meio de diferentes famílias de aprendizado; no espaço de modalidades, pelas diferentes combinações de fontes de informação; e no espaço de decisão, considerando a complementaridade de diferentes erros produzidos pelos classificadores \cite{britto2014dynamic, monteiro2021diversity}.

Inicialmente, os 15 candidatos com melhor desempenho médio em F1-macro foram selecionados para a etapa de formação final da \textit{pool}. Em seguida, foi adotada uma estratégia gulosa iterativa para selecionar 7 especialistas, com base no critério definido pela Equação \ref{eq:score_candidato}.

\begin{equation}
\label{eq:score_candidato}
\text{Score}(C_i) = \alpha \cdot P(C_i) + (1 - \alpha) \cdot D(C_i, S)
\end{equation}
em que:
\begin{itemize}
    \item $P(C_i)$ representa o desempenho normalizado do candidato ($C_i$), obtido durante a validação cruzada;
    \item $D(C_i, S)$ representa seu \textit{disagreement} médio em relação ao conjunto ($S$) de classificadores já selecionados na \textit{pool};
    \item $\alpha$ controla a importância relativa entre desempenho e diversidade, definido como 0.7 no experimento.
\end{itemize}



\subsection{Seleção Dinâmica de Classificadores}
A DS parte da premissa de que nenhum classificador da \textit{pool} é competente em todo o espaço de entrada, embora cada um possa sê-lo em regiões específicas~\cite{britto2014dynamic, cruz2018dynamic}. Para cada instância de teste, a competência dos membros é estimada em uma região de competência, usualmente formada pelos $k$ vizinhos mais próximos em um conjunto próprio, o DSEL, mantido separado do treino porque os classificadores tenderiam a acertar as instâncias que já viram. Com base nessa estimativa, seleciona-se um único classificador (\textit{Dynamic Classifier Selection}) ou um subconjunto deles (\textit{Dynamic Ensemble Selection}).

A competência dos sete especialistas, já treinados no TRAIN e mantidos congelados, é estimada no DSEL. Para cada instância de teste, buscam-se os $k = 7$ vizinhos mais próximos no DSEL, em um espaço que concatena as quatro modalidades, cada uma padronizada (ajuste no TRAIN) e normalizada pela norma L2, para que tenham o mesmo peso na distância.

% *********************************************
Foram avaliadas três estratégias de DS: OLA, KNORA-E e KNORA-U. No OLA (\textit{Overall Local Accuracy})~\cite{woods1997combination}, decide sozinho o especialista com maior acurácia entre os vizinhos. Já no KNORA-E (\textit{K-Nearest Oracles Eliminate})~\cite{ko2008dynamic}, votam os especialistas que acertaram todos os $k$ vizinhos, com redução da vizinhança quando nenhum o faz. Enquanto no KNORA-U (\textit{K-Nearest Oracles Union})~\cite{ko2008dynamic}, todos votam, com peso igual ao número de vizinhos classificados corretamente.

A incongruência entra apenas na região de competência. Sem roteamento, os vizinhos são buscados no DSEL inteiro; com roteamento, apenas entre as instâncias do mesmo grupo (A/B ou conflito) da instância de teste. Como restringir a vizinhança altera a estimativa independentemente do critério, cada roteamento foi comparado também a um controle aleatório, com grupos sorteados do mesmo tamanho dos reais (média de 10 sorteios). Um ganho atribuível à incongruência deve, portanto, aparecer contra os dois.


\subsection{Mixture of Experts}

% explicar o que é, conceitos
Arquiteturas \textit{Mixture of Experts} (Mistura de Especialistas, em português) combinam as previsões de múltiplos especialistas por meio de um mecanismo de roteamento (\textit{gate}), que aprende a atribuir diferentes pesos a cada especialista de acordo com a entrada apresentada \cite{jacobs1991adaptive}.

Neste trabalho, foi implementada uma arquitetura \textit{Soft-MoE}, que utiliza a \textit{pool} já previamente formada, composta por 7 especialistas cujas previsões são mantidas fixas durante o treinamento do mecanismo de combinação. Nessa variante, todos os especialistas contribuem para a predição final, com pesos determinados dinamicamente pelo mecanismo de roteamento.

O \textit{gate} foi implementado como uma rede neural do tipo \textit{Multi-Layer Perceptron} (MLP), que recebe como entrada a concatenação das representações das quatro modalidades (texto, contexto, áudio e rosto), totalizando 2320 dimensões. Sua arquitetura consiste em uma camada oculta com 32 neurônios, ativação ReLU e \textit{Dropout} de 0.2, seguida de uma camada de saída com sete neurônios. A função \textit{Softmax} converte as saídas em pesos cuja soma é igual a um. Esses pesos são utilizados para combinar as probabilidades de sarcasmo produzidas pelos especialistas, gerando a predição final.

O treinamento do \textit{gate} é realizado exclusivamente no conjunto DSEL, utilizando entropia cruzada binária (BCE) como função de perda e o otimizador Adam, com taxa de aprendizado 0.001, \textit{weight decay} de 0.0001 e lotes de 32 amostras. São permitidas até 200 épocas, com parada antecipada após 20 épocas sem melhora.

% finalizar explicitando que a intenção é utilizar como comparação com a nossa proposta
% A comparação estrutural planejada contrapõe um \textit{gate} MoE padrão, aprendido ponta a ponta sobre as representações multimodais.
Diferentemente da seleção dinâmica de classificadores, que estima a competência local dos especialistas, o MoE aprende diretamente os pesos de combinação a partir das representações multimodais. Neste estudo, essa abordagem é utilizada como referência comparativa para avaliar o desempenho da combinação aprendida em relação aos métodos de seleção dinâmica propostos.

% ~\cite{yu2024mmoe}.
% cite{jacobs1991adaptive}
%mecanismos de \textit{gate} que ponderam dinamicamente entre especialistas para capturar incongruência semântica entre texto e imagem~\cite{zhao2025beyond, jana2025midre};.

\subsection{Protocolo Experimental e Avaliação Estatística}

O protocolo utilizou 10 \textit{splits} com partição independente de falante, de modo a evitar o risco de enviesamento dos modelos \cite{ray2022multimodal, wang2026bottleneck}. Em cada repetição, aproximadamente 70\% das instâncias foram destinadas ao treinamento (TRAIN), 15\% à seleção dinâmica (DSEL) e 15\% ao teste (TEST), gerados por \texttt{GroupShuffleSplit} agrupado pela identidade do falante, garantindo que nenhum indivíduo estivesse presente em mais de um conjunto. Devido à distribuição desigual das falas entre os 27 falantes do MUStARD++, foram aceitos apenas \textit{splits} com 120 a 240 instâncias em DSEL e TEST e proporção de sarcasmo entre 40\% e 60\%. A Tabela \ref{tab:dados_splits} apresenta a distribuição final.

\input{tables/datasplit}

O TRAIN foi utilizado para seleção, ajuste de hiperparâmetros (validação cruzada estratificada de 5 \textit{folds} agrupados por falante) e treinamento dos classificadores; o DSEL, para estimar a competência na DS e treinar o \textit{gate} do Soft-MoE; e o TEST, exclusivamente para a avaliação final. 

A configuração principal da Seleção Dinâmica utilizou $k=7$ vizinhos e uma \textit{pool} de sete especialistas, formada com peso de desempenho de 0.7 e diversidade 0.3. Foram avaliados OLA, KNORA-E e KNORA-U sem roteamento e com a região de competência restrita pelos grupos de incongruência. Os controles aleatórios do roteamento usaram 10 sorteios por split.

Como abordagens comparativas, empregaram-se o Soft-MoE e classificadores estáticos (SVM, LR e RF) treinados com a concatenação das quatro modalidades — texto, contexto, áudio e visual —, com os melhores hiperparâmetros obtidos previamente por validação cruzada no conjunto TRAIN na etapa de formação da \textit{pool}, além do voto majoritário entre eles. Por aplicarem a mesma regra a todas as instâncias, esses modelos permitem avaliar se a adaptação por instância promovida pela DS oferece vantagem em relação à fusão multimodal estática. Além disso, foram utilizados o melhor e o pior membro da \textit{pool} (selecionados pelo desempenho no DSEL), o voto majoritário da \textit{pool} e um oráculo, que representa o teto de desempenho da \textit{pool}, considerando uma instância correta sempre que pelo menos um dos sete especialistas produz a classificação correta.

O F1-macro foi a métrica principal, por ponderar igualmente o desempenho nas duas classes; acurácia, precisão e \textit{recall} foram complementares. Para cada método, os resultados foram agregados entre os 10 \textit{splits}, reportando-se média, desvio-padrão e os desempenhos mínimo e máximo.

Na comparação principal, aplicou-se o teste de Friedman~\cite{friedman1937use} ao F1-macro de 11 métodos: os quatro estáticos, o Soft-MoE, o melhor e o pior membro da \textit{pool}, o voto da \textit{pool} e as três estratégias de DS sem roteamento, considerando cada \textit{split} como um bloco. Seguiu-se o pós-teste de Nemenyi para comparações entre pares de métodos~\cite{demsar2006statistical}, com nível de significância $\alpha=0.05$. Os roteamentos ficaram fora desse teste por serem avaliados separadamente, e o oráculo, reportado apenas em acurácia, também não foi incluído.

Para avaliar especificamente o efeito do roteamento por incongruência, foram realizadas quatro comparações pareadas: (i) roteamento A/B e ausência de roteamento; (ii) roteamento A/B e seu controle aleatório; (iii) roteamento por conflito e ausência de roteamento; e (iv) roteamento por conflito e seu respectivo controle aleatório. Realizadas para OLA, KNORA-E e KNORA-U (12 testes), usaram o teste de Wilcoxon pareado~\cite{wilcoxon1945individual} sobre o F1-macro dos 10 \textit{splits}, com correção de Holm~\cite{holm1979simple} ($\alpha=0.05$). A sensibilidade ao número de vizinhos ($k \in \{3,5,7,11,15\}$) foi avaliada de forma descritiva.

Por fim, foi medido o custo computacional das principais etapas (tempo de execução, tempo aproximado de CPU e pico de memória), para avaliar também a viabilidade prática dos métodos.

\section{Resultados e Discussão}
\label{sec:results}

Os resultados por \textit{split} estão no repositório (Apêndice \ref{apendice}).

\subsection{Comparação principal}

A Tabela \ref{tab:main} contempla as comparações principais do experimento com a média nos 10 \textit{splits}.

\input{tables/results main}

Os classificadores estáticos com voto dos 3 modelos e SVM obtiveram o maior F1-macro médio (80.1), seguidos pelo KNORA-U sem roteamento e pelo voto da \textit{pool} (79.2) e pelo Soft-MoE (78.3). O teste de Friedman, aplicado aos 11 métodos nos 10 \textit{splits}, rejeitou a hipótese de equivalência ($\chi^2 = 21.32$, $p = 0.019$), mas nenhum par superou a distância crítica de Nemenyi (4.77): a maior diferença de rank médio foi de 4.55, entre o voto dos três estáticos e o pior membro da \textit{pool}.

Dessa forma, a seleção dinâmica não superou os classificadores estáticos e o Soft-MoE empatou com OLA e KNORA-E (78.2) e ficou abaixo do KNORA-U, sem diferença significativa. A seleção por competência local mostrou-se, portanto, comparável à combinação aprendida pelo \textit{gate}, sem vantagem clara de nenhuma das duas. A variação entre \textit{splits} (desvio-padrão de 5 a 7 p.p.) superou as diferenças entre métodos.

\subsection{Efeito do roteamento por incongruência}

Cada estratégia de DS foi executada com os cinco roteamentos: sem roteamento, A/B, conflito e os controles aleatórios de cada um. A Tabela~\ref{tab:roteamento} mostra que as diferenças ficaram abaixo de um ponto e que os controles aleatórios reproduziram os valores dos roteamentos reais. A maior diferença ocorreu no KNORA-E com roteamento A/B, 0.5 ponto acima da versão sem roteamento.

\input{tables/ablation roteamento}

Para verificar se essas diferenças são sistemáticas, cada roteamento foi comparado à versão sem roteamento e ao seu controle, por estratégia, com o teste de Wilcoxon pareado sobre os 10 \textit{splits}. Os 12 testes foram corrigidos por Holm, para controlar o erro de múltiplas comparações, e nenhum foi significativo (\textit{p} ajustado = 1.0 em todos). Como o roteamento real não se distinguiu de grupos sorteados de mesmo tamanho, o resultado sugere que ele atuou apenas como restrição da vizinhança, sem acrescentar informação sobre a competência dos especialistas. O pequeno benefício do KNORA-U também ocorre sem roteamento, indicando que vem da combinação dos especialistas e não da incongruência. O resultado também não depende do tamanho da vizinhança: entre $k = 3$ e $k = 15$, o F1-macro médio de cada estratégia e roteamento variou no máximo 0.7 ponto.

\subsection{Por que o roteamento não ajudou?}
\label{subsec:ablacao}

Foram investigadas duas causas possíveis: um sinal de incongruência pouco confiável e uma \textit{pool} pouco diversa. Em cada análise, alterou-se um fator por vez em relação à configuração principal, que foi fixada antes da avaliação no teste; a única exceção é a calibração dos extratores, que combina dois ajustes.

\subsubsection{Os grupos medem de fato a incongruência emocional?}

Se os extratores de emoção errarem com frequência, os grupos A/B e de conflito não representam o desacordo real entre canais e o roteamento não teria como funcionar. Para verificar essa possibilidade, a polaridade estimada em cada modalidade foi comparada com a emoção explícita anotada no MUStARD++, convertida para as mesmas três polaridades. Essa anotação descreve a fala com todas as modalidades e não é uma referência exata para cada canal, mas, se os extratores capturassem a emoção da fala, deveriam concordar razoavelmente com ela.

A concordância foi baixa em todas as modalidades (Tabela~\ref{tab:diagnostico}), com kappa de no máximo 0.27 no texto. O contexto é o caso menos comparável, já que descreve falas anteriores e não a fala anotada. O principal problema, porém, foi o excesso de saídas neutras: os extratores classificaram como neutras cerca de metade das falas em que os anotadores identificaram emoção. Isso afeta diretamente os grupos. Das 173 instâncias do grupo A (no conjunto completo; nos conjuntos de teste dos splits, o grupo tem de 16 a 49 falas), 141 (81\%) estão nele porque as quatro modalidades foram classificadas como neutras, e não porque expressam a mesma emoção. Na prática, o agrupamento A/B separa falas sem emoção detectada das demais, e não falas congruentes de incongruentes, o que oferece uma explicação para os resultados obtidos.

\input{tables/ablation diagnostico}

Para tentar melhorar esse sinal, foram testados dois ajustes, escolhidos apenas com o TRAIN: (i) a polaridade passou a ser definida por um limiar $\tau$ sobre a valência (positiva se $V_m > \tau$, negativa se $V_m < -\tau$ e neutra entre os dois), com escolha de $\tau$ por modalidade e, no texto e no contexto, também do extrator; e (ii) o rosto foi reextraído considerando apenas o personagem que fala, com o ViT \texttt{trpakov/vit-face-expression}~\cite{trpakov_vitface}. Fora do TRAIN, nenhum ajuste melhorou a concordância de forma consistente: o texto subiu no DSEL (kappa de 0.263 para 0.294), mas não no TEST (de 0.292 para 0.286), as demais modalidades caíram, e o rosto do falante concordou tanto quanto o original (kappa em torno de 0.1). Com os grupos recalculados, o roteamento também continuou sem diferença significativa. A limitação parece estar, portanto, na capacidade dos extratores disponíveis de reconhecer emoção nessas falas, e não na regra usada para definir a polaridade.

\subsubsection{Uma \textit{pool} mais diversa favoreceria a seleção dinâmica?}

O oráculo atingiu 91.0 de acurácia, contra em torno de 79.0 dos métodos dinâmicos: cerca de 12\% das falas seriam acertadas por algum membro, mas o seletor escolheu errado. O erro da \textit{pool} concentrou-se nas falas sem sarcasmo (15.6\%, contra 2.9\% nas sarcásticas), e a \textit{pool} apresentou pouca diversidade (\textit{kappa} entre 0.67 e 0.83 por \textit{split}): 98.6\% dos seus membros usam o contexto (Figura \ref{fig:pool}), e especialistas parecidos erram nas mesmas instâncias, o que deixa pouco espaço para a seleção escolher entre eles. Foram avaliadas três variantes na formação da \textit{pool}: com $\alpha$ 0.2 e 0, no lugar de 0.7 (aumentando o peso de diversidade para 0.8 e 1.0) e um limite de quatro membros com contexto. Para não usar o teste na escolha, as variantes foram comparadas pelo F1-macro do voto da \textit{pool} no DSEL.

% Figura de exemplo
\begin{figure}[h] 
    \centering 
    \includegraphics[width=0.67\linewidth]{images/resultados-pool.png} 
    \caption{Porcentagem dos 70 membros (7 por split) que usam cada modalidade ou classificador.} 
    \label{fig:pool} 
\end{figure}

O ganho de diversidade foi real, mas não se converteu em desempenho. O oráculo subiu de 91.0\% para até 93.2\%, enquanto o KNORA-U variou no máximo 0.4 ponto, e nenhuma variante diferiu da configuração padrão (Wilcoxon com Holm, \textit{p} ajustado = 1.0). Os membros que entraram sem contexto são individualmente mais fracos, o que sugere que os novos acertos do oráculo vêm de classificadores que a estimativa de competência local raramente escolhe.

\subsection{Custo computacional}

A Tabela \ref{tab:custo-computacional} resume o custo das etapas. A etapa mais cara foi a escolha da \textit{pool} (cerca de 24.5 minutos por \textit{split}), enquanto a DS completa, com todos os roteamentos, levou cerca de 106 segundos, e a decisão da DS custou cerca de 0.05 milissegundos por fala. O Soft-MoE acrescenta 74503 parâmetros treináveis e teve o maior pico de memória entre os métodos de classificação (1239 MB).

\input{tables/custo-computacional}

\section{Conclusões}
\label{sec:conclusions}

Este trabalho propôs aplicar estratégias de Seleção Dinâmica de classificadores à detecção de sarcasmo, usando a incongruência intermodal como critério de competência. Foram realizadas duas análises principais: (1) como a DS se compara aos classificadores estáticos e ao Soft-MoE; e (2) se a incongruência, como critério explícito de competência, melhora o desempenho da DS.

Nos 10 \textit{splits} por falante, o melhor F1-macro médio foi dos classificadores estáticos SVM e voto dos três modelos (80.1). A DS ficou de 0.8 a 0.9 p.p. abaixo deles e empatou com o Soft-MoE: o KNORA-U (79.2) ficou cerca de 0.9 p.p. acima do Soft-MoE (78.3), sem diferença significativa no pós-teste de Nemenyi, e o mesmo F1 ocorre sem roteamento. O roteamento por incongruência não alterou o resultado de forma significativa. Assim, nenhuma hipótese do trabalho foi confirmada estatisticamente.

Os resultados obtidos têm limitações. O MUStARD++ possui apenas 27 falantes, muito desiguais, e, mesmo com o critério de aceitação, o \textit{split} pesa mais que o método; além disso, com 10 \textit{splits}, os testes têm pouco poder para diferenças de 1 a 2 pontos, e ``não significativo'' não implica equivalência. O sinal de incongruência também é ruidoso: a polaridade estimada concorda pouco com a anotação da base (\textit{kappa} de no máximo 0.27, no texto), e os ajustes testados não a melhoraram de forma consistente fora do treino, o que pode ter limitado o efeito do roteamento.

Por outro lado, o oráculo atingiu 91.0 de acurácia, e aumentar a diversidade da \textit{pool} elevou esse teto a até 93.2 sem ganho para a DS, o que indica potencial, mas o seletor atual não o aproveita. Como trabalhos futuros, recomenda-se melhorar o seletor, com alternativas à extração de emoções, e também explorar a incongruência como entrada do \textit{gate} do MoE.

\appendices
\section{Repositório}
\label{apendice}
O código-fonte e os arquivos desenvolvidos neste trabalho estão disponíveis no repositório do GitHub:

\begin{center}
\url{https://github.com/berkembrockisabella/sarcasm-dynamic-selection}
\end{center}

\bibliographystyle{IEEEtran}
\bibliography{bibtex}

\end{document}